\begin{solapartat}
\punts{0,2} $m_m v r_m = \dfrac{h}{2\pi} \Rightarrow v = \dfrac{h}{2\pi}\,\dfrac{1}{m_m r_m}$

\punts{0,25} D'altra banda, considerant que el muó descriu un moviment circular uniforme al voltant del nucli, la seva acceleració centrípeta és: $a = \dfrac{v^2}{r_m} = \left(\dfrac{h}{2\pi}\right)^2\dfrac{1}{m_m^2 r_m^3}$

\punts{0,1} La segona llei de Newton estableix que la força normal és: $F = m_m a = \left(\dfrac{h}{2\pi}\right)^2\dfrac{1}{m_m r_m^3}$

\punts{0,1} i la llei de Coulomb: $F = k\,\dfrac{e^2}{r_m^2}$

\punts{0,25} Per tant obtenim que $k\,\dfrac{e^2}{r_m^2} = \left(\dfrac{h}{2\pi}\right)^2\dfrac{1}{m_m r_m^3}$

$r_m = \left(\dfrac{h}{2\pi}\right)^2\dfrac{1}{m_m k e^2} = 2,65\times 10^{-13}\ \mathrm{m}$ \quad \punts{0,1}

\punts{0,25} De l'anterior expressió podem comprovar que el radi de l'orbita és inversament proporcional a la massa del leptó, per tant el radi de l'electró serà 200 vegades més gran que el radi del muó. Llavors, serà molt més gran l'hidrogen que l'hidrogen muònic.
\end{solapartat}

\begin{solapartat}
\punts{0,5} L'energia del fotó emès és igual a la energia perduda per l'electró quan salta de l'orbita de més energia a la fonamental, és a dir,:
\[ E_{\textit{fotó}} = -(E_1 - E_2) = 3,266\times 10^{-16}\ \mathrm{J} \]

\punts{0,25} Per altre banda l'energia del fotó és: $E_{\textit{fotó}} = hf = h\,\dfrac{c}{\lambda}$

I finalment,

\punts{0,25} $f = \dfrac{E_{\textit{fotó}}}{h} = 4,93\times 10^{17}\ \mathrm{Hz}$

\punts{0,25} $\lambda = h\,\dfrac{c}{E_{\textit{fotó}}} = \dfrac{c}{f} = 6,09\times 10^{-10}\ \mathrm{m}$
\end{solapartat}
